% Commutative diagram from the long-exact-sequence chapter.
% Source: book/tex/homology/long-exact.tex (lines 592-599).
\documentclass[border=2mm]{standalone}
\usepackage{amssymb,amsmath}
\usepackage{tikz-cd}
\newcommand{\ZZ}{\mathbb{Z}}
\newcommand{\eps}{\varepsilon}
\newcommand{\id}{\mathrm{id}}
\begin{document}
	\begin{tikzcd}
		\cdots \ar[r]
			& \underbrace{H_2(U) \oplus H_2(V)}_{= 0} \ar[r]
			& \underbrace{H_2(X)}_{\cong \ZZ} \ar[lld, "\partial"', swap] \\
		\underbrace{H_1(U \cap V)}_{\cong \ZZ \oplus \ZZ} \ar[r, "\phi"']
			& \underbrace{H_1(U) \oplus H_1(V)}_{\cong \ZZ \oplus \ZZ} \ar[r]
			& \cdots
	\end{tikzcd}
\end{document}
